\documentclass[12pt, a4paper]{article}

\usepackage[utf8]{inputenc}
\usepackage[spanish]{babel}
\usepackage{geometry}
\usepackage{booktabs}
\usepackage{longtable}
\usepackage{array}
\usepackage{xcolor}
\usepackage{colortbl}
\usepackage{titlesec}
\usepackage{hyperref}
\usepackage{enumitem}

\geometry{margin=2.5cm}

% Colores
\definecolor{cabecera}{RGB}{30, 90, 160}
\definecolor{filasombreada}{RGB}{220, 232, 250}

% Estilo de sección
\titleformat{\section}{\large\bfseries\color{cabecera}}{}{0em}{}[\titlerule]

% --------------------------------------------------
\begin{document}

\begin{center}
    {\LARGE \textbf{Reparto de Tareas}}\\[0.4em]
    {\large \textit{Space Invaders — Proyecto de Programación}}\\[0.3em]
    {\small Fecha: \today}
\end{center}

\vspace{1em}
\hrule
\vspace{1.5em}

% --------------------------------------------------
\section{Miembros del equipo}

\begin{itemize}[leftmargin=2em]
    \item \textbf{Miembro 1:} Nombre Apellido
    \item \textbf{Miembro 2:} Nombre Apellido
    \item \textbf{Miembro 3:} Nombre Apellido
    % Añadir o eliminar miembros según corresponda
\end{itemize}

\vspace{1.5em}

% --------------------------------------------------
\section{Tabla de reparto de tareas}

\begin{longtable}{>{\bfseries}p{0.04\textwidth} p{0.28\textwidth} p{0.40\textwidth} p{0.16\textwidth}}
    \rowcolor{cabecera}
    \textcolor{white}{\#} &
    \textcolor{white}{Tarea} &
    \textcolor{white}{Descripción} &
    \textcolor{white}{Responsable(s)} \\
    \midrule
    \endfirsthead

    \rowcolor{cabecera}
    \textcolor{white}{\#} &
    \textcolor{white}{Tarea} &
    \textcolor{white}{Descripción} &
    \textcolor{white}{Responsable(s)} \\
    \midrule
    \endhead

    % ---- Fila de ejemplo (duplicar para añadir más tareas) ----
    \rowcolor{filasombreada}
    1 & Diseño del modelo (MVC) &
    Definición de las clases del modelo: \texttt{Espacio}, \texttt{Jugador}, \texttt{Enemigo}, \texttt{Naves}, \texttt{Disparo}. Implementación de la lógica del juego. &
    Miembro 1 \\
    \midrule

    2 & Patrón Observer &
    Implementación de las interfaces \texttt{Observable} y \texttt{Observador}. Conexión entre modelo y vista. &
    Miembro 2 \\
    \midrule

    \rowcolor{filasombreada}
    3 & Vista y controlador &
    Desarrollo de \texttt{MainFrame} y \texttt{StartFrame}. Pintado del tablero, captura de eventos de teclado y actualización de la interfaz. &
    Miembro 3 \\
    \midrule

    4 & Lógica de disparo y colisiones &
    Detección de colisiones entre disparos y naves enemigas/jugador. Control del ciclo de vida de los disparos. &
    Miembro 1, Miembro 2 \\
    \midrule

    \rowcolor{filasombreada}
    5 & Movimiento de enemigos &
    Algoritmo de movimiento grupal de los enemigos, cambio de dirección al llegar a los bordes y descenso de filas. &
    Miembro 2 \\
    \midrule

    6 & Gestión del juego (Main / bucle) &
    Configuración del \texttt{Main}, hilo del juego, control de estados (inicio, partida, fin) y reinicio. &
    Miembro 3 \\
    \midrule

    \rowcolor{filasombreada}
    7 & Documentación &
    Redacción de la memoria, diagramas UML, justificación del patrón MVC+Observer y este documento de reparto. &
    Todos \\
    \midrule

    % ---- Añadir nuevas tareas aquí ----
    8 & & & \\

    \bottomrule
\end{longtable}

\vspace{1.5em}

% --------------------------------------------------
\section{Detalle individual de contribuciones}

% ---- Repetir este bloque para cada miembro ----
\subsection*{Miembro 1 — Nombre Apellido}
\begin{itemize}[leftmargin=2em]
    \item \textbf{Tarea 1 — Diseño del modelo:} descripción breve de lo que hizo concretamente (clases creadas, métodos implementados, decisiones de diseño tomadas\ldots).
    \item \textbf{Tarea 4 — Lógica de disparo y colisiones:} descripción breve.
\end{itemize}

\subsection*{Miembro 2 — Nombre Apellido}
\begin{itemize}[leftmargin=2em]
    \item \textbf{Tarea 2 — Patrón Observer:} descripción breve.
    \item \textbf{Tarea 4 — Lógica de disparo y colisiones:} descripción breve.
    \item \textbf{Tarea 5 — Movimiento de enemigos:} descripción breve.
\end{itemize}

\subsection*{Miembro 3 — Nombre Apellido}
\begin{itemize}[leftmargin=2em]
    \item \textbf{Tarea 3 — Vista y controlador:} descripción breve.
    \item \textbf{Tarea 6 — Gestión del juego:} descripción breve.
\end{itemize}

\vspace{1.5em}

% --------------------------------------------------
\section{Observaciones generales}

\begin{itemize}[leftmargin=2em]
    \item Escribir aquí cualquier nota sobre la coordinación del equipo, herramientas usadas (Git, etc.) o dificultades encontradas.
\end{itemize}

\end{document}
